%% ==============================
\chapter{\iflanguage{ngerman}{Zusammenfassung und Ausblick}{Conclusion}}
\label{sec:conclusion}
%% ==============================

\section{Zusammenfassung}
\label{sec:conclusion-zusammenfassung}

Ausgangspunkt dieser Arbeit war eine Lücke, die sich aus der Literatur nicht schließen ließ. Für das dezentrale Gesundheitsdatenmanagement liegen zahlreiche blockchainbasierte Architekturvorschläge vor, hashgraphbasierte Ansätze sind dagegen kaum untersucht, und die publizierten Evaluationsergebnisse sind wegen unterschiedlicher Messgrößen, Lastprofile und Systemgrenzen untereinander nicht vergleichbar (Kapitel~\ref{cha:state_of_the_art}). Ein belastbarer Vergleich beider Ledger-Familien setzte daher eine eigene Messung unter identischen Bedingungen voraus.

Als Gegenstand dieser Messung diente die Referenzarchitektur von Erler et al., in der die Gesundheitsdaten selbst außerhalb des Ledgers verbleiben und die Ledger-Schicht ausschließlich als Vertrauensanker für Bezeichner, Schemas, Credential Definitions und Widerrufsinformationen dient. Genau diese Abgrenzung machte den Austausch der Ledger-Schicht überhaupt zu einem kontrollierten Experiment: Bleibt die Agenten- und Kommunikationsschicht unverändert, sind gemessene Unterschiede der ausgetauschten Schicht zuzuordnen.

Der Austausch selbst erwies sich als Wechsel der Systemklasse und nicht als Wechsel der Implementierung. Hyperledger Indy führt den Zustand der registrierten Identitätsobjekte selbst und setzt Schreibberechtigungen über ein Rollenmodell unmittelbar im Ledger durch; der Hedera Consensus Service ordnet Nachrichten innerhalb eines Topics manipulationsresistent, interpretiert ihren Inhalt jedoch nicht. Alles, was die eine Variante im Ledger erbringt, musste auf der anderen oberhalb des Ledgers rekonstruiert werden. Kapitel~\ref{cha:konzept} hat diese Rekonstruktion in sieben Entwurfsentscheidungen zerlegt und jede davon gegen die zuvor abgeleiteten Anforderungen abgewogen. Zwei dieser Entscheidungen bestimmen das Ergebnis der Arbeit maßgeblich: die Ableitung des Objektzustands aus der Nachrichtenhistorie bei jedem Lesezugriff (E5, Abschnitt~\ref{subsec:e5-leseweg}) und die Führung der Vertrauenshierarchie über ein eigenes Topic mit nachgelagerter Prüfung beim Lesen (E6, Abschnitt~\ref{subsec:e6-vertrauenshierarchie}).

Kapitel~\ref{cha:implementierung} hat dieses Konzept als Proof-of-Concept umgesetzt und den Messaufbau so eingerichtet, dass beide Varianten unter derselben Grundlast, mit derselben Kennzahlendefinition und mit derselben Wiederholungszahl vermessen werden konnten. Fünf Lastszenarien -- Basismessung, Laststeigerung, Datenwachstum, Knotenausfall und Zugriffskontrolle -- wurden auf beiden Varianten durchgeführt und in Kapitel~\ref{sec:results} berichtet. Innerhalb desselben Kapitels wurden diese Ergebnisse entlang der Frage eingeordnet, welche der beobachteten Unterschiede aus der Architektur der beiden Distributed Ledger folgen und welche lediglich aus dem Softwarestand, mit dem sie angesprochen wurden (Abschnitt~\ref{sec:results-einordnung}).

\section{Beantwortung der Forschungsfragen}
\label{sec:conclusion-forschungsfragen}

\paragraph{FF1 -- Performanz.} Zwischen beiden Varianten besteht ein erheblicher Abstand in der operationsspezifischen Latenz. Schreibvorgänge liegen auf der HCS-Variante zwischen 4,3\,s und 13,0\,s, auf der Indy-Variante zwischen 0,1\,s und rund 3,0\,s, mit einer zweigipfligen Ausnahme bis etwa 6,0\,s. Beim Lesen fällt der Abstand noch deutlicher aus: Für diejenigen Objekttypen, die den in Abschnitt~\ref{sec:impl-befunde} nachgewiesenen Zeitboden treffen, stehen rund 5\,s gegenüber wenigen Hundertstelsekunden. Da eine einzelne Präsentationsprüfung mehrfach lesend auf den Ledger zugreift, ist dies der praktisch bedeutsamste Einzelbefund der Arbeit.

Diese Zahlen sind jedoch nur eingeschränkt als Technologieaussage lesbar. Ein wesentlicher Teil der HCS-Werte gruppiert sich so eng um einen festen Wert, dass nicht das Konsensverfahren, sondern ein von der Client-Bibliothek abgewartetes Zeitfenster die Dauer bestimmt; der kausale Nachweis gelang über die Variation dieses Fensters und wurde durch die Objekttypselektivität des Bodens unabhängig bestätigt. Die Werte beschreiben damit, was ein Aufrufer mit diesem Softwarestand erhält, und nicht, wie lange das Netzwerk benötigt. Der dritte in FF1 genannte Aspekt, der Ressourcenverbrauch, bleibt unbeantwortet: Er wurde auf keiner der beiden Varianten erhoben (Abschnitt~\ref{sec:results-ressourcen}). Diese Teilfrage ist damit offen und wird nicht durch eine qualitative Ersatzaussage überdeckt.

\paragraph{FF2 -- Skalierbarkeit.} Bei steigender Nebenläufigkeit verhalten sich beide Varianten grundlegend verschieden, und zwar nicht graduell, sondern in der Verlaufsform. Die Indy-Variante wächst über elf Laststufen bis zu einem Durchsatzplateau bei 55 parallelen Nutzern, bleibt dabei durchgehend fehlerfrei und wird langsamer, bevor sie an ihre Grenze stößt. Die HCS-Variante bricht dagegen zwischen fünf und zehn parallelen Nutzern unvermittelt ein; die Fehlerrate springt auf über 98\,\%, während die Antwortzeit an derselben Stelle fällt, weil scheiternde Anfragen sofort beantwortet werden. Ein allein auf die Antwortzeit gestütztes Sättigungskriterium hätte diesen Zusammenbruch als Verbesserung gelesen -- ein methodisches Nebenergebnis, das für künftige Vergleiche dieser Art von Bedeutung ist.

Die Ursache des Zusammenbruchs ließ sich bis auf die Quellcodeebene zurückverfolgen: Der Konsensknoten meldet unter Last eine spezifikationsgemäß vorübergehende Überlastung, die die verwendete Client-Bibliothek nicht wiederholt, sondern unverändert als endgültigen Fehler durchreicht. Auch dieser Befund ist damit der Software und nicht dem Konsensverfahren zuzuschreiben.

Bei wachsendem Objektbestand ergibt sich ein symmetrisches Bild: Auf keiner der beiden Varianten war ein Effekt auf die Lesezeit auflösbar, auf der Indy-Variante bis 2700 Objekte, auf der HCS-Variante bis 300. Dieses Ergebnis ist ausdrücklich als obere Schranke und nicht als Nullbefund formuliert. Der Entwurf selbst sagt vorher, dass der Effekt erst bei hinreichend großen Beständen sichtbar wird; die Messung zeigt daher, wo er nicht liegt, und nicht, dass es ihn nicht gibt. Aufschlussreicher als das ausbleibende Ergebnis war der Grund für den Abbruch der HCS-Bestandserweiterung: Eine mehrstündige Schreiblast genügte, um die Lesbarkeit des Netzwerks zum Erliegen zu bringen, ohne dass ein Knotenausfall nötig gewesen wäre.

\paragraph{FF3 -- Sicherheit.} Die Unterschiede in der Durchsetzbarkeit der Schreibautorisierung sind kategorialer und nicht gradueller Natur. Die HCS-Variante bindet den Schreibzugriff an einen Schlüssel je Objekt, die Indy-Variante an eine netzwerkweit geführte Rolle je Akteur. Beide Netzwerke weisen unautorisierte Schreibversuche zurück; die Indy-Variante gestattet darüber hinaus eine abgestufte Autorisierung, die sowohl zwischen Rollen als auch zwischen Operationstypen wirkt und deren Ablehnungen reicher beobachtbar sind. Das Modell der HCS-Variante ist deshalb nicht schwächer, sondern andersartig, und die fehlende anwendungsseitige Delegationsstufe ist ein nicht umgesetzter, kein protokollseitiger Nachteil.

Eine Einschränkung begrenzt die Reichweite dieser Antwort erheblich und betrifft beide Varianten gleichermaßen: Die Agentenschicht signiert jede Ledger-Transaktion mit ihrer eigenen Kennung, unabhängig davon, welcher Aussteller in der Anfrage angegeben ist. Über den Zugriffsweg, den eine Anwendung im Betrieb tatsächlich nimmt, unterscheidet die Kontrolle damit nicht mehr zwischen Ausstellern, sondern nur noch zwischen Agenten. Die in Abschnitt~\ref{sec:results-s5} nachgewiesenen Unterschiede sind Eigenschaften der Ledger-Schicht und beschreiben nicht, was ein Betreiber ohne weitere Maßnahmen tatsächlich an Zugriffskontrolle erhält.

Zur Konsens-Fehlertoleranz fällt die Antwort schmaler aus als die Frage. Auf der HCS-Variante wurde nicht die Abwesenheit von Fehlertoleranz gemessen, sondern die Unfähigkeit des amtierenden Softwarestands, sie zugänglich zu machen: Die Adressierung ist strukturell auf einen einzigen Konsensknoten festgelegt, sodass ein vierknotiges Netzwerk keinen realisierten Mehrwert gegenüber einem einzelnen Knoten liefert, obwohl die Redundanz betrieben wird. Ob das Hashgraph-Verfahren selbst fehlertolerant ist, bleibt durch diese Arbeit unbeantwortet. Auf der Indy-Variante löste sich der Ausfall eines konsensführenden Knotens nach etwa einer Minute über einen protokollinternen Sichtwechsel ohne äußeren Eingriff auf, während der Störung selbst keine Anfrage beantwortet wurde; ob dies auf eine Beschlussunfähigkeit des Netzwerks oder auf eine fehlende Ausweichlogik der Client-Bibliothek zurückgeht, ließ sich nicht klären.

\paragraph{Hauptforschungsfrage.} Beide Technologien eignen sich als Verifiable Data Registry einer SSI-basierten Architektur für dezentrales Gesundheitsdatenmanagement, jedoch in unterschiedlichem Reifegrad und mit unterschiedlichem Integrationsaufwand. Die funktionale Eignung des Hedera Consensus Service ist belegt: Sämtliche Kernfunktionen der Ledger-Schicht ließen sich mit HCS-Mechanismen nachbilden und in allen fünf Szenarien tatsächlich ausführen; kein Szenario scheiterte daran, dass die Technologie eine benötigte Fähigkeit nicht besäße. Der Preis dieser Substitution ist architektonischer Natur und im Entwurf benannt: Zustandshaltung und Berechtigungsprüfung wandern aus dem Ledger in die darüberliegende Schicht.

Die gemessene Eignung fällt zum Erhebungszeitpunkt dennoch klar zugunsten der Indy-Variante aus -- in der Latenz, im Sättigungsverhalten und in der nutzbaren Fehlertoleranz. Entscheidend für die Interpretation ist, dass diese Nachteile überwiegend nicht dem Hashgraph-Konsensverfahren zuzuschreiben sind, sondern dem Stand der Client-Bibliotheken, mit denen es angesprochen wurde. Drei der vier Hauptbefunde -- der Zeitboden, der Lastzusammenbruch und der nicht nutzbare Knotenausfall -- folgen diesem Muster und ließen sich jeweils bis auf die Bibliotheksebene zurückverfolgen. Diese Trennung ist keine Entlastung der HCS-Variante: Wer sie heute betreiben will, betreibt sie mit eben diesen Bibliotheken. Sie ist eine Aussage darüber, was sich ändern kann und was nicht. Für eine Entscheidung im Jahr der Messung folgt daraus die Indy-Variante; für eine Bewertung der Technologie folgt daraus, dass der Vergleich mit einem neueren Softwarestand zu wiederholen wäre, bevor dem Hashgraph-Verfahren die gemessenen Nachteile zugeschrieben werden.

\section{Fazit}
\label{sec:conclusion-fazit}

Über alle Szenarien hinweg lässt sich die funktionale Substitution als gelungen bezeichnen. Sämtliche Kernfunktionen der Ledger-Schicht -- Registrierung von Bezeichnern, Schema, Credential Definition, Widerruf und Zugriffskontrolle -- konnten mit HCS-Mechanismen nachgebildet und in den Szenarien S1 bis S5 tatsächlich ausgeführt werden. Die in Kapitel~\ref{cha:konzept} entworfene Zielarchitektur ist in diesem Sinne tragfähig. Kein Szenario scheiterte daran, dass die Technologie eine benötigte Fähigkeit nicht besäße.

Die gemessenen Nachteile sind überwiegend Eigenschaften der eingesetzten Client-Bibliotheken und nicht des Hashgraph-Konsensverfahrens. Dieses Muster trägt drei der vier Hauptbefunde -- den Zeitboden in S1, den Zusammenbruch in S2 und den nicht messbaren Knotenausfall in S4 -- und rechtfertigt die durchgehend getrennte Lesart von Technologie und Softwarestand. Diese Trennung ist keine Entlastung der HCS-Variante: Wer sie heute betreiben will, betreibt sie mit eben diesen Bibliotheken. Sie ist eine Aussage darüber, was sich ändern kann und was nicht.

Der architektonische Preis der in E5 gewählten Zustandsableitung beim Lesen ist real, im untersuchten Rahmen jedoch nicht dramatisch. Er tritt allerdings nicht dort auf, wo er zuerst vermutet wurde: Ein Effekt der Bestandsgröße war nicht auflösbar, während die eigentliche Belastung sich in der Betriebsfestigkeit unter anhaltender Schreiblast zeigte. Der Speicherbedarf, der den unmittelbarsten Preis der Historienführung darstellt, blieb ungemessen.

Sicherheitsseitig ist das Rollenmodell der Indy-Variante reicher und in seinen Ablehnungen besser beobachtbar. Das Modell der HCS-Variante ist deshalb nicht schwächer, sondern andersartig: objektgebundene Schlüsselkontrolle statt Rollenhierarchie. Die fehlende anwendungsseitige Delegationsstufe ist ein nicht umgesetzter, kein protokollseitiger Nachteil. Beide Modelle teilen jedoch eine Schwäche, die keinem von beiden anzulasten ist: Über den im Betrieb üblichen Zugriffsweg werden sie von der Agentenschicht gleichermaßen unterlaufen.

Schließlich ist die Evidenzstärke der einzelnen Befunde bewusst unterschiedlich und wird durchgehend so ausgewiesen. Manches ist kausal bewiesen, etwa der Zeitboden über die Variation des Zeitfensters und der Zusammenbruch über den Quellcode des Schreibpfads. Anderes ist beobachtet, aber nicht mechanistisch erklärt, etwa die Zweigipfligkeit auf der Indy-Variante und die Ursache der Nichtverfügbarkeit bei Knotenausfall. Wieder anderes ist als obere Schranke und nicht als Nullbefund formuliert, etwa das Ergebnis zum Datenwachstum. Diese Abstufung wird nicht zu einer einheitlichen Sicherheitsstufe verflacht, weil die Aussagekraft der Arbeit gerade darin besteht, dass die Unterschiede zwischen diesen Graden benannt bleiben.

Der Beitrag dieser Arbeit lässt sich daraus in drei Punkten zusammenfassen. Der inhaltliche Beitrag liegt in der erstmaligen Konzeption und Umsetzung einer SSI-basierten Verifiable Data Registry auf einem reinen Ordnungsdienst im Gesundheitskontext, einschließlich der vollständigen Dokumentation der sieben dafür nötigen Entwurfsentscheidungen und ihrer Abwägung gegen die Anforderungen. Die Zielarchitektur liegt damit nicht als Vorschlag, sondern als vermessener Prototyp vor.

Der methodische Beitrag liegt in einem Messaufbau, der beide Varianten unter identischen Bedingungen vergleichbar macht, und in der eben beschriebenen durchgehenden Trennung zwischen Aussagen über eine Technologie und Aussagen über ihren Softwarestand. Diese Trennung wurde nicht nachträglich behauptet, sondern in jedem Einzelfall über die Quellcodeebene oder ein kausales Experiment belegt.

Als Nebenergebnis sind vier zuvor nicht gemeldete Einschränkungen der verwendeten Client-Bibliotheken entstanden, die sich in einem systematischen Reifegradgefälle zwischen Lese- und Schreibpfad derselben Bibliotheksgeneration zusammenfassen lassen.

\section{Ausblick}
\label{sec:conclusion-ausblick}

Aus den benannten Grenzen ergeben sich mehrere Anschlussarbeiten, die sich in ihrem Aufwand deutlich unterscheiden.

Am unmittelbarsten anschlussfähig ist die Wiederholung der Szenarien S2 und S4 mit einer Client-Bibliothek, die bei vorübergehender Ablehnung wiederholt und mehrere Konsensknoten anspricht. Für wenigstens eine der vier beschriebenen Einschränkungen ist belegt, dass neuere Bibliotheksversionen sie beheben; für die übrigen ist es zu prüfen. Ein solcher Durchlauf ließe sich als eigens gekennzeichneter Vergleichsarm führen und würde unmittelbar zeigen, welcher Teil der gemessenen Nachteile mit dem Softwarestand entfällt. Im Rahmen dieser Arbeit wurde darauf verzichtet, weil ein Wechsel des Softwarestands mitten in der Messreihe die Vergleichbarkeit aller zuvor erhobenen Werte aufgehoben hätte.

Die Erhebung des Ressourcenverbrauchs nach NFA3 und des Speicherbedarfs nach NFA5 ist die Lücke mit dem besten Verhältnis von Aufwand zu Erkenntnisgewinn. Sie erfordert keine neue Implementierung, sondern lediglich eine zusätzliche Beobachtung während der ohnehin durchgeführten Szenarien -- und sie würde zugleich klären, ob der beobachtete Abbruch der HCS-Bestandserweiterung eine Frage der Ressourcenzuteilung oder eine strukturelle Eigenschaft ist.

Die Bestandsserie zum Datenwachstum sollte auf der HCS-Variante auf das auf der Indy-Variante erreichte Niveau gebracht werden; Voraussetzung dafür ist eine Ressourcenausstattung der Mirror Node, die eine mehrstündige Schreiblast übersteht. Ergänzend wäre die zweite, bislang nicht untersuchte Wachstumsrichtung interessant: die Lesezeit eines einzelnen Topics bei wachsender eigener Nachrichtenhistorie, die auf der HCS-Variante der eigentliche Prüfstein der Entscheidung E5 wäre.

Die anwendungsseitige Prüfstufe nach E6 ist eine eigenständige Implementierungsarbeit. Sie würde nicht nur das Szenario S5 auf der HCS-Variante vervollständigen, sondern ist zusammen mit Befund B4 auch die Antwort auf die Frage, wie eine Zugriffskontrolle aussehen müsste, die nicht von der Agentenschicht unterlaufen wird -- eine Frage, die beide Varianten gleichermaßen betrifft und über den Rahmen dieser Arbeit hinausweist.

Schließlich sollte der Mechanismus hinter der Nichtverfügbarkeit bei Knotenausfall auf der Indy-Variante geklärt werden. Dafür genügt vermutlich eine Wiederholung des Szenarios unter Aufbewahrung der Knotenprotokolle aller vier Knoten -- ein geringer Mehraufwand, der die in dieser Arbeit offengebliebene Unterscheidung zwischen Beschlussunfähigkeit und fehlender Ausweichlogik entscheiden würde.